\chapter{Normalização e Apresentação Gráfica}

A normalização assegura unidade visual, facilita a leitura e permite que os trabalhos sejam avaliados por critérios comuns. No padrão da UnDF, as decisões de apresentação consideram as normas ABNT aplicáveis e escolhas institucionais para os pontos em que há margem de definição.

\section{Papel, margens e espaçamento}

O trabalho deve ser apresentado em papel A4. O padrão institucional utiliza margem superior e esquerda de 3 cm, margem inferior e direita de 2 cm. O corpo textual utiliza espaçamento de uma linha e meia, com exceções para elementos que exigem espaçamento simples, como citações longas, notas, referências, natureza do trabalho e fontes de ilustrações.

\requisitonormativo

O formato A4 e as margens indicadas seguem os requisitos de apresentação da ABNT NBR 14724:2024. O espaçamento de uma linha e meia no corpo acompanha a recomendação normativa e foi adotado pelo projeto. O recuo de parágrafo e a ausência de espaço adicional são decisões institucionais validadas no âmbito do projeto.

\recomendacaonormativa

\decisaoinstitucional

\section{Fonte e legibilidade}

A ABNT recomenda tamanho 12 para o texto e tamanho menor e uniforme para exceções, mas não determina uma família tipográfica. Para compatibilizar os modelos editáveis, o projeto adota o seguinte padrão institucional:

\decisaoinstitucional

\begin{longtable}{p{0.42\linewidth}p{0.21\linewidth}p{0.25\linewidth}}
\toprule
Elemento & Tamanho & Tratamento \\
\midrule
Corpo do texto & 12 pt & Família serifada institucional \\
Seção primária & 12 pt & Negrito e caixa alta \\
Demais seções & 12 pt & Negrito e diferenciação hierárquica consistente \\
Citação direta longa & 10 pt & Recuo de 4 cm, espaço simples, sem aspas \\
Notas de rodapé & 10 pt & Espaço simples \\
Legendas e fontes & 10 pt & Espaço simples \\
Natureza do trabalho & 10 pt & Espaço simples \\
Conteúdo de quadros e tabelas & 10 pt ou 12 pt & Uniforme e legível no documento \\
\bottomrule
\end{longtable}

\fluxopendente

A família tipográfica institucional ainda precisa de deliberação. Até que exista padrão aprovado, o estudante deve seguir a orientação vigente do curso e manter uma única família legível e uniforme em todo o trabalho, sem combinar fontes apenas por efeito visual.

Mais importante do que escolher uma fonte incomum é manter consistência. Misturar muitas fontes, tamanhos e estilos prejudica a aparência institucional do trabalho.

\section{Parágrafos}

O padrão institucional adota recuo de 1,25 cm nos parágrafos, inclusive no primeiro parágrafo após títulos. O texto principal não deve usar espaçamento extra entre parágrafos como substituto do recuo, salvo em elementos específicos definidos pelo modelo.

\section{Paginação}

A capa não apresenta número visível. A folha de rosto é contada, mas a numeração visível aparece a partir da parte textual. O número deve permanecer em posição estável e seguir a contagem até os elementos pós-textuais.

A página dos dados internacionais de catalogação não entra na contagem. No anverso, o número aparece no canto superior direito, a 2 cm das bordas superior e direita. O Apêndice \ref{ap:paginacao} mostra a sequência completa.

\section{Títulos e numeração progressiva}

Os títulos devem refletir a hierarquia do trabalho. Seções primárias, secundárias, terciárias, quaternárias e quinárias devem ser numeradas de forma progressiva, respeitando a organização do conteúdo.

O padrão institucional atual apresenta as seções primárias no formato ``1 INTRODUÇÃO'', sem a palavra ``Capítulo'', com número e título na mesma linha e título em caixa alta.

\decisaoinstitucional

O uso de caixa alta na seção primária e os destaques gráficos dos demais níveis são escolhas da UnDF. Pela ABNT NBR 6024:2012, a numeração deve limitar-se até a seção quinária e não se usa ponto, hífen, travessão ou outro sinal entre o indicativo e o título.

Títulos sem indicativo numérico, como resumo, listas, sumário, referências, glossário, apêndices, anexos e índice, recebem tratamento próprio. Elementos pré-textuais não entram no sumário; referências e demais elementos pós-textuais entram sem indicativo numérico.

\section{Alíneas e subalíneas}

Alíneas são usadas quando o texto precisa apresentar enumeração interna sem transformar cada item em seção. Devem manter paralelismo gramatical: se o primeiro item começa com verbo, os demais também devem seguir estrutura semelhante; se começa com substantivo, os demais devem manter o padrão.

Use alíneas com moderação. Listas muito longas podem indicar que o conteúdo deveria ser reescrito em parágrafos, quadro ou tabela.

\section{Siglas}

Na primeira ocorrência, a sigla deve aparecer após a forma por extenso. Depois disso, a sigla pode ser usada sozinha.

Exemplo: Universidade do Distrito Federal Professor Jorge Amaury Maia Nunes (UnDF). Nas ocorrências seguintes, pode-se usar UnDF.

Siglas pouco conhecidas, siglas numerosas ou siglas específicas de área podem justificar uma lista de abreviaturas e siglas.

\section{Abreviaturas e meses}

Abreviaturas devem seguir uso padronizado e não devem prejudicar a clareza. Em referências, datas e meses precisam respeitar o padrão normativo adotado. O estudante deve evitar misturar meses por extenso, abreviados e em outro idioma na mesma lista de referências.

\section{Impressão e versão digital}

Mesmo quando a entrega é digital, a apresentação deve considerar leitura em tela e eventual impressão. Figuras, tabelas, gráficos e cores devem permanecer legíveis. Evite depender exclusivamente de cor para transmitir informação, pois a impressão em preto e branco pode prejudicar a compreensão.

\section{Sumário}

O sumário deve reproduzir os títulos e a paginação do trabalho. Não deve incluir elementos pré-textuais. Deve ser conferido depois da versão final, pois alterações de página e título podem desatualizar suas entradas.

\section{Consistência}

A normalização exige consistência. Se uma figura recebe fonte abaixo, todas devem receber. Se a banca usa determinado formato, todos os membros devem seguir o mesmo padrão. Se as referências usam um padrão de data de acesso, ele deve ser mantido.

\caixaatencao{A revisão da versão de entrega deve ser feita no PDF, não apenas no arquivo editável. Alguns problemas de quebra de página, paginação, figuras e sumário só aparecem claramente no documento gerado ou exportado.}
